% Source: https://artofproblemsolving.com/community/c3h2516735_a_couple_of_problems
% Stable identity: aops-topic-2516735-q06
\textbf{Day 1:} $49.30\cdot1.01=49.793$, which rounds to $50$.

\textbf{Later days.} Suppose that at the start of a day the amount owed is a whole number $n$ with $50\le n\le148$. Then $1.01n=n+\dfrac{n}{100}$.
\begin{itemize}
\item If $50\le n\le99$: $0.5\le\dfrac n{100}<1$, so $n+\dfrac n{100}$ rounds to $n+1$ (for $n=50$ the value $50.5$ is rounded up).
\item If $100\le n\le148$: $1.01n=(n+1)+\dfrac{n-100}{100}$ with $0\le\dfrac{n-100}{100}<0.5$, which again rounds to $n+1$.
\end{itemize}
So after day 1 the debt grows by exactly $\$1$ per day: after $d$ days it is $49+d$ dollars (for $d\le100$). After $30$ days the amount owed is $\boxed{\$79}$, of which
\[79-49.30=\boxed{\$29.70}\]
is interest. (Without rounding, the debt would be $49.30\cdot1.01^{30}\approx\$66.45$.)
